\section{The prayers and readings of the Sunday}

\subsection{A man who has put his trust in someone}

The singer tells God that he has lifted his soul to him and that he trusts him. He has still to find out what his enemies will see when they look at him. If he had put his confidence in someone powerless to help him, they would have reason to laugh. Bellarmine hears this exposure to ridicule in the psalm. The suppliant has rooted his expectation so publicly in God that to be let down would shame him before others. Yet even before he asks for a way out, he can recall what happens to others who have waited for this same Lord. The turn from one man's enemies to all those who wait gives a congregation room to sing in the singular.\footnote{Ps 24(25):1--4. Robert Bellarmine, \work{Commentary on the Book of Psalms}, Ps 24, exposition of vv. 1--4, in John O'Sullivan's abridged English.}

The verse between the two singings of the entrance antiphon asks to be shown a path. The man who dreads ridicule has a reason to ask this. Another man's willingness to believe in him, or his own success in an undertaking, would hardly tell him whether he had taken the right road. Bellarmine understands the paths as the way to live so as to please God. A petition that began in the hearing of enemies becomes a request God alone can grant. When the antiphon returns, the singer has also given glory to Father, Son, and Holy Spirit. The short doxology remains at this point in the Mass even though the longer Gloria after the Kyrie is omitted at a Mass of the Advent season.\footnote{\work{Missale Romanum} (Vatican, 1962), p. [1], no. 1 and following rubric.}

\subsection{The peril that follows one's own sins}

The Collect begins where the prayer of Psalm 79 begins to grow urgent. Israel has invoked God as the shepherd who guides Joseph like a sheep. The people call upon him to stir up his strength and come to save them. They will have to turn again to God, and more than once in the psalm they ask him to turn them. They ask to see his face shine on them. In Augustine's explanation, one can turn one's face away from God of one's own will and yet have to ask for God's help to turn back. The shepherd whom the people want to see comes among them in the flesh.\footnote{Ps 79:2--4 (80:1--3). Augustine, \work{Enarrationes in Psalmos}, on Ps 79, §§2--3, historical English under Psalm LXXX.}

The Sunday prayer tells us something further about the people who ask for rescue. The peril approaching them comes from their own sins. The worshipper who first heard about enemies in the Introit must now reckon with what he himself has done. The priest uses the first person plural. He asks that the Lord may protect those who are threatened, deliver them, and bring them to safety. In the Latin, \emph{te protegente} and \emph{te liberante} mean that the protection and the releasing are the Lord's acts. Even the request to be found deserving of rescue takes that form. The congregation answers Amen to a prayer addressed to the Son himself. The close calls him God and speaks of his life with the Father in the unity of the Holy Spirit. The identity of the One being asked to come is thus already explicit before either lesson.\footnote{1962 Missal, p. [1], no. 2. Lasance and Walsh, \work{The New Roman Missal}, First Sunday of Advent, p. 103, with its reference to Prayer Conclusions, p. 819.}

\subsection{What Paul says to people who already believe}

In the preceding sentences Paul has told the Romans that the debt to love the other person endures. He has named several commandments that a person who loves his neighbor will keep. The biblical sentence at Rom 13:11 begins with a conjunction pointing back to that instruction. Then Paul speaks of knowing the time. At Mass the priest reads an address to the brethren in place of the connective, and the Latin omits the noun for time, \emph{tempus}. The word for the hour, \emph{hora}, remains. Both historical English texts contain a time or season at the opening. A reader of Chrysostom and Aquinas will find those commentators explaining the longer biblical opening. The comparison in the next sentence belongs to the Bible and the Mass alike. The people in question already know what it is to believe. Since then something has happened to the distance between them and salvation.\footnote{Rom 13:8--14. Compare the Clementine Vulgate, Rom 13:11, with the 1962 Missal, no. 3, and Pustet's \work{Missale Romanum} (1862), First Sunday of Advent, p. 2.}

It takes Paul three images to say what a waking Christian does. A sleeper can throw off what he wore in bed. Someone taking the field puts on armor. Someone walking abroad in the daytime goes where others can see him. Each image corrects an easy misunderstanding of the last one. To get out of bed still leaves a man to decide how to spend the day. A man prepared for combat might think his business was to defeat another person. Paul counts quarrels and jealousy among the things he is to put away. At the point where one might expect to learn which virtues the armor stands for, Paul gives his hearers the name of a person to put on. The Mass ends the lesson there, at the name of the Lord Jesus Christ. The biblical sentence goes on to warn about making provision for the flesh and its desires.\footnote{Rom 13:12--14. John Chrysostom, \work{Homilies on Romans}, Homily 24, on vv. 12--14.}

\subsection{Two requests between the lessons}

The Gradual begins with the assurance about those who wait. The threatened man, his enemies and his fear of ridicule have dropped out. A vocative at the end of the Latin response turns the assurance directly to the Lord. Laurence Shepherd renders this extra address in his English of the Gradual. In the English of Lasance and Walsh one has to reach the verse to hear the first ``O Lord.'' The verse again asks to know the Lord's paths. The request also varies slightly in Latin. The verb in the Introit asks God to \emph{show} his ways. Here the singers ask him to \emph{make them known}. The singers pass over the clause of the biblical psalm that asks for shame for those who act unjustly. What the congregation takes up from the psalm after hearing Paul name particular sins is a promise to the people who hope and a petition for instruction.\footnote{Ps 24(25):3--4. 1962 Missal, nos. 1, 4. Prosper Guéranger, \work{The Liturgical Year}, Advent, trans. Laurence Shepherd, First Sunday of Advent, p. 134.}

In the Alleluia the singers ask both to \emph{see} mercy and to \emph{receive} salvation. In the psalm as a whole Israel looks back on God's favor to his land, on pardon and the turning away of anger. The people can remember having been helped and still need him to turn them again. In the biblical wording a change of tense marks the difference. What God has done gives people grounds for what they are about to ask. Augustine hears the past tense at the beginning as the speech of a prophet certain enough of God's coming gift to speak of it as accomplished. He distinguishes that prophetic use from the reality of Israel's earlier captivities and releases. When he reaches the petition sung in the Alleluia he asks what salvation God could show them. His answer is Christ. A people who ask to be shown mercy may therefore be asking to see the face of the person who will come to them.\footnote{Ps 84:2--8 (85:1--7). Augustine, \work{Enarrationes}, on Ps 84, §§2--4, 6 (English Psalm LXXXV).}

\subsection{What Jesus asks the disciples to see}

The people listening to Jesus in Luke have admired the temple and heard him predict its destruction. Earlier in the discourse he warned his disciples about people who would announce that the time had come and solicit followers; he told them to refuse. He has spoken about Jerusalem surrounded by armies before the Mass takes up his words about the whole inhabited world. Now sky, sea and earth are all involved. People are already suffering from what they expect, while the event itself still lies ahead. The disciples share their surroundings. From Jesus they learn that they can raise their heads at the beginning of these things. The reason is a personal promise. It is \emph{your} redemption that is coming near.\footnote{Lk 21:5--8, 20--28. For the setting of the discourse see \hyperref[sec:date-location]{Scriptural Date and Location}.}

Then Jesus directs the disciples' eyes to a fig tree and to all the trees. In the Missal they are producing fruit. People accustomed to watching them know from experience that summer is close. Jesus invites the disciples to exercise an ordinary capacity for recognition under conditions in which fear could make them unable to use it. The one who knows summer is near cannot thereby give the day and hour of the first fruit's ripening. He can, however, say something true about what lies ahead. In the explanation Jesus speaks of the kingdom of God. In his assurance about the future he speaks of his own words. The listeners have moved from seeing the familiar signs to depending on what he has told them about events outside their experience.\footnote{Lk 21:29--33, 1962 Missal, p. 2, no. 6.}

\subsection{The returning voice and the offered gifts}

The Latin of the Offertory is exactly the antiphon with which the Mass opened. The psalm verse, the doxology and the repetition belong to the entrance. In Lasance and Walsh a small feature of the English would otherwise obscure the identity of the two Latin antiphons. Their Introit adds the address ``O Lord'' after the first ``Thee,'' as the Douay Bible does. Their Offertory begins simply ``To Thee have I lifted up my soul.'' In both places the Missal begins with the same address and the same completed act of lifting. At the altar, people who are offering something of their own make the same declaration of dependence that they made on entering.\footnote{Ps 24(25):1--3. 1962 Missal, nos. 1, 7. Lasance and Walsh, First Sunday of Advent, pp. 103, 105.}

The prayer over the gifts was traditionally called the Secret. Its Latin has \emph{haec sacra}, the holy things here before the priest, as the subject. The people for whom he prays are the ones to undergo a change. He asks that the gifts may cleanse the people and bring them, purer, to the very source of these holy things. There is a difficulty at the last word. \emph{Principium} can signify a beginning or an origin. Both the Lasance and the Shepherd English understand a personal origin here. In the printed English the gift has an author, before whom the cleansed can come. There are thus two reasons to speak of approaching in this prayer. The people physically draw near to receive, and they ask to be made fit for the One they approach.\footnote{1962 Missal, p. 2, no. 8. Lasance and Walsh, First Sunday of Advent, p. 105. Guéranger, \work{Advent}, First Sunday of Advent, pp. 136--137.}

\subsection{To whom the Church gives thanks}

At the opening of the Preface the priest calls for the lifting of hearts. The people answer for themselves before he calls them to give thanks. The priest then addresses the Father and speaks to him about his only-begotten Son and the Holy Spirit. He says what the Church believes about the Father by revelation and what follows about the others. The glory confessed of the Father belongs equally to the Son and the Spirit. The conclusion distinguishes what belongs to each person, the one essence they share, and the equality of their majesty. A reader might wonder how the Son can be sent to creatures in need if he has the same majesty as the Father. Augustine takes up this difficulty in his inquiry about divine missions. Being sent concerns the Son's relation to the one who sends him. One cannot infer that he is inferior in divinity from the fact that he is sent.\footnote{1962 Missal, p. 2 after no. 8 and p. 293, no. 1082. Augustine, \work{De Trinitate}, III, preface §3, trans. Arthur West Haddan.}

When the priest has spoken of the persons and their common majesty, he names those who already praise God. Angels, archangels, cherubim and seraphim are distinguished in the list, and their praise has one voice. The congregation can join that praise in the Sanctus. Those who have just heard that the powers of heaven will be shaken are allowed to say to the Lord of hosts that heaven and earth are full of his glory. In Luke the shaking of heaven belongs to the prospect of judgment. The adoration is happening at the altar at the time the priest says these words.

\subsection{Receiving and still asking to receive}

The singers return to Psalm 84 at Communion. Much has happened in the psalm between the Alleluia petition and this verse. The speaker has undertaken to hear what God will say to his people. Mercy and truth have met, justice and peace have kissed, and truth has sprung from earth while justice looked down from heaven. God gives goodness, and the earth gives its fruit. The psalm lets the earth be truly productive while naming divine goodness first. Both verbs in Latin are the same future, \emph{dabit}, ``will give.'' A promise can be sung by people already eating because the sacrament and the visible harvest are different kinds of gift. What one can mean by the earth's fruit changes according to what one takes God to have given.\footnote{Ps 84:8--13 (85:7--12). 1962 Missal, p. 2, no. 9.}

The Postcommunion immediately offers another instance of people still asking for a gift. Psalm 47 celebrates the security of God's city and remembers mercy received in his temple. In Latin the psalm says \emph{suscepimus}, ``we have received.'' The prayer says \emph{suscipiamus}, ``may we receive.'' Augustine's exposition of the psalm has ``in the midst of thy people'' where the Missal has ``in the midst of thy temple.'' He can speak directly of the assembly on the strength of the words he has before him. Here it is the situation at Mass that lets us hear the prayer as the prayer of an assembled people. Having received the sacrament, they ask to receive mercy in the temple so that they may prepare with due reverence for the feast of their redemption. Where Lasance speaks of the coming festival, the Latin has a plural, the approaching solemnities.\footnote{Ps 47:10 (48:9). Augustine, \work{Enarrationes}, on Ps 47, §§8--9 (English Psalm XLVIII). 1962 Missal, p. 2, no. 10.}
